% ===========================================================================
%  PART 1 --- Introduction, the BST problem, the AVL idea, height, proof,
%             Fibonacci property.
%  Follows the requested flow:
%    intro -> what is a BST -> what problem -> AVL idea -> how AVL solves it
%    -> height property -> proof -> build-up -> golden ratio -> exact count
%  Wording rule: short sentences, common words, no idioms.
%
%  Height convention throughout the deck: h counts EDGES, so a single node is
%  height 0 and N(0)=1. A lecture that calls a single node "height 1" gets
%  N(h) = F(h+2)-1; ours is F(h+3)-1. Same fact, shifted index. Do not mix them.
% ===========================================================================

% ------------------------------------------------------------------ title
\begin{frame}[plain,label=p1]
  % 5 px three-hue top bar: the three parts, in the order they arrive.
  % Anchored to the page itself, not to the text block, so it bleeds edge to
  % edge; the same picture carries the progress mark at the opposite corner.
  \begin{tikzpicture}[remember picture,overlay]
    \fill[avlpone]   (current page.north west) rectangle
                     ($(current page.north west)+(0.3333\paperwidth,-5\pxl)$);
    \fill[avlptwo]   ($(current page.north west)+(0.3333\paperwidth,0)$) rectangle
                     ($(current page.north west)+(0.6666\paperwidth,-5\pxl)$);
    \fill[avlpthree] ($(current page.north west)+(0.6666\paperwidth,0)$) rectangle
                     ($(current page.north east)+(0,-5\pxl)$);
    \fill[avlpone]   (current page.south west) rectangle
                     ($(current page.south west)+
                       (\dimexpr\paperwidth/\inserttotalframenumber\relax,3\pxl)$);
    % Course mark sits on the footer line the content frames use: same 84 px
    % gutter, clear of the progress tick in the corner below it.
    \node[anchor=west,inner sep=0pt]
      at ($(current page.south west)+(84\pxl,42\pxl)$) {\eyebrow{CSE 200}};
  \end{tikzpicture}
  \vfill
  \begin{center}
    {\avlfDeck\wxbold\color{avlink}AVL Trees}\\[18\pxl]
    {\avlfLead\wmed\color{avlfib}Trees which balance all on their own}\\[34\pxl]
    % Fibonacci glyph, flanked by hairlines: the thread, planted before a word
    % of it is spoken. leafSep 13 px, levSep 16 px, per the spec.
    \begin{tikzpicture}[x=1\pxl,y=1\pxl,
        avldot/.style={circle,draw=avlfib!62!avlground,
                       fill=avlfib!62!avlground,inner sep=0pt,minimum size=6.4\pxl},
        avledge/.style={draw=avlfib!42!avlground,line width=1.4\pxl}]
      \setavlscale{13}{16}
      \resetavl
      \drawfib{7}{0}{-46}
      \draw[avlrule,line width=1\pxl] (-243,-46) -- (-155,-46);
      \draw[avlrule,line width=1\pxl] ( 155,-46) -- ( 243,-46);
    \end{tikzpicture}\\[30\pxl]
    % Three equal-width columns, so the roll centres under its own name and the
    % blocks space evenly however long the names are. The roll takes the
    % eyebrow treatment -- faint and tracked -- so it reads as a subtitle to the
    % name rather than competing with it.
    \begin{minipage}[b]{0.30\textwidth}\centering
      {\avlfBody\wmed\color{avlbody}Md Ishtiak Moin}\\[7\pxl]
      {\avlfEyebrow\wsemi\color{avlfaint}\avltrack{2305011}}
    \end{minipage}%
    \begin{minipage}[b]{0.30\textwidth}\centering
      {\avlfBody\wmed\color{avlbody}Sheikh Raufur Rahim}\\[7\pxl]
      {\avlfEyebrow\wsemi\color{avlfaint}\avltrack{2305010}}
    \end{minipage}%
    \begin{minipage}[b]{0.30\textwidth}\centering
      {\avlfBody\wmed\color{avlbody}Abhijnan Podder}\\[7\pxl]
      {\avlfEyebrow\wsemi\color{avlfaint}\avltrack{2305007}}
    \end{minipage}
  \end{center}
  \vfill
\end{frame}

\avldivider{1}{The problem}{When can a search tree stop being fast?}{avlpone}{avlponeA}{avlponeB}{avlponeC}

\setavlpart{1}{Part 1 --- The problem}{avlpone}

% ------------------------------------------------------------------ what is a BST
\begin{frame}[label=p1]{First, the rule every BST follows}
  \vspace{0.3em}
  \begin{columns}[T,onlytextwidth]
    \begin{column}{0.52\textwidth}
      {\avlfLead\wbold\color{avlink}Binary Search Tree}\\[0.7em]
      {\color{avlbody}For \emph{every} node:}
      \begin{itemize}
        \setlength{\itemsep}{0.35em}
        \item everything on the \fib{left} is \fib{smaller}
        \item everything on the \fib{right} is \fib{bigger}
      \end{itemize}
      \vspace{0.6em}
      {\avlfSec\color{avlmuted}To search, you compare once and drop one whole
       side.}\\[0.9em]
      \begin{avlpanel*}
        Cost of a search $=$ \fib{the height of the tree}.
      \end{avlpanel*}
    \end{column}
    \begin{column}{0.44\textwidth}
      \begin{center}
      \begin{tikzpicture}[x=1cm,y=1cm,scale=0.98]
        \node[avlkey] (n4) at (0,0)       {4};
        \node[avlkey] (n2) at (-1.35,-1.2) {2};
        \node[avlkey] (n6) at (1.35,-1.2)  {6};
        \node[avlkey] (n1) at (-2.1,-2.4)  {1};
        \node[avlkey] (n3) at (-0.6,-2.4)  {3};
        \node[avlkey] (n5) at (0.6,-2.4)   {5};
        \node[avlkey] (n7) at (2.1,-2.4)   {7};
        \foreach \a/\b in {n4/n2,n4/n6,n2/n1,n2/n3,n6/n5,n6/n7}
          \draw[avledge] (\a) -- (\b);
        \node[avlannot] at (-2.05,-0.50) {smaller};
        \node[avlannot] at ( 2.05,-0.50) {bigger};
        % annotation row under the figure: hairline, key, hairline
        \draw[avlrule,line width=1\pxl] (-2.55,-3.15) -- (-0.95,-3.15);
        \draw[avlrule,line width=1\pxl] ( 0.95,-3.15) -- ( 2.55,-3.15);
        \node[color=avlfib,font=\avlfSec\wmed] at (0,-3.15) {height $=2$};
      \end{tikzpicture}
      \end{center}
    \end{column}
  \end{columns}

  \note{Say the rule out loud. Left is smaller, right is bigger. That is all a
  BST is. And because you throw away half the tree at each step, the cost of a
  search is just the height.}
\end{frame}

% ------------------------------------------------------------------ the problem
\begin{frame}[label=p1]{The problem: sorted input destroys it}
  % Captions sit ABOVE each tree, a hairline divides the two halves, and the
  % colour carries the verdict: green is the legal-and-short tree, red the list.
  % The green nodes on the left are the search path to 7 -- three of them.
  \vspace{0.1em}
  \begin{center}
  \begin{tikzpicture}[x=1cm,y=1cm,scale=0.94,every node/.style={transform shape}]
    \figdivider{0}{0.95}{-3.75}

    \begin{scope}[xshift=-3.55cm,yshift=-0.55cm]
      \node[avlannot,text=avlmuted] at (0,1.35) {inserted 4, 2, 6, 1, 3, 5, 7};
      \node[avlkeyok]  (b4) at (0,0)        {4};
      \node[avlkey]    (b2) at (-1.25,-1.0) {2};
      \node[avlkeyok]  (b6) at (1.25,-1.0)  {6};
      \node[avlkey]    (b1) at (-1.95,-2.0) {1};
      \node[avlkey]    (b3) at (-0.55,-2.0) {3};
      \node[avlkey]    (b5) at (0.55,-2.0)  {5};
      \node[avlkeyok]  (b7) at (1.95,-2.0)  {7};
      \foreach \a/\b in {b4/b2,b2/b1,b2/b3,b6/b5}
        \draw[avledge] (\a) -- (\b);
      \foreach \a/\b in {b4/b6,b6/b7}
        \draw[draw=avlbalance,line width=1.6\pxl] (\a) -- (\b);
      % both verdicts share one baseline: this scope sits 0.40 lower than the
      % right-hand one, so its label y is 0.40 higher to compensate
      \node[color=avlbalance,font=\avlfStmt\wbold] at (0,-2.95) {3 steps};
    \end{scope}

    \begin{scope}[xshift=3.55cm,yshift=-0.15cm]
      \node[avlannot,text=avlmuted] at (0,0.95) {inserted 1, 2, 3, 4, 5, 6, 7};
      % step must exceed the disc diameter or the edges vanish inside the nodes
      \foreach \k [count=\i from 0] in {1,2,3,4,5,6,7} {
        \pgfmathsetmacro{\yy}{-\i*0.44}
        \pgfmathsetmacro{\xx}{-1.20+\i*0.40}
        \node[avlkeyalert,minimum size=34\pxl,font=\avlfAnnot] (c\k) at (\xx,\yy) {\k};
      }
      \foreach \a/\b in {c1/c2,c2/c3,c3/c4,c4/c5,c5/c6,c6/c7}
        \draw[draw=avlalert,line width=1.6\pxl] (\a) -- (\b);
      \node[color=avlalert,font=\avlfStmt\wbold] at (0,-3.35) {7 steps};
    \end{scope}
  \end{tikzpicture}
  \end{center}

  \vspace{0.4em}
  \begin{center}
    {\avlfSec\color{avlmuted}Same seven keys. Both are legal BSTs.
     The right one is now a \viol{list}, not a tree.}
  \end{center}

  \note{Both trees are legal. The right one is what you get when the data
  arrives already sorted --- IDs, dates, imported records. This is not a rare
  case. It is the most common case. Search cost goes from log n to n.}
\end{frame}

% ------------------------------------------------------------------ the idea
\begin{frame}[label=p1]{The idea: check the height at every node}
  \vspace{0.6em}
  \begin{center}
    {\avlfLead\color{avlbody}The tree got tall because one side grew and
     \emph{nobody checked}.}
  \end{center}
  \vspace{0.6em}
  \begin{center}
  \begin{tikzpicture}[x=1cm,y=1cm,scale=0.90,every node/.style={transform shape}]
    \node[avlkey,minimum size=42\pxl] (r) at (0,0) {};
    % subtrees as ember triangles, each carrying its height as the symbol
    \node[avltri,minimum size=2.3cm,font=\avlfLead,text=avlfib] (L) at (-1.75,-1.55) {$a$};
    \node[avltri,minimum size=1.75cm,font=\avlfLead,text=avlfib] (R) at (1.75,-1.25) {$b$};
    \draw[avledge] (r) -- (L);
    \draw[avledge] (r) -- (R);
    % height rulers, one per side, measured from the root row to each base
    \heightbrace{-3.05}{-0.30}{-2.32}{height $a$}{avlfaint}
    \heightbrace{3.05}{-0.30}{-1.75}{height $b$}{avlfaint}
    \draw[avlrule,line width=1\pxl,dotted] (-3.05,-0.30) -- (3.05,-0.30);
    \draw[avlrule,line width=1\pxl] (-2.55,-2.72) -- (2.55,-2.72);
    \node[avlannot,anchor=north] at (0,-2.80) {one node, two subtree heights};
  \end{tikzpicture}
  \end{center}

  \vspace{0.5em}
  \begin{center}
    \begin{avlpanel}[0.62\textwidth]
      The rule: $a$ and $b$ may differ by at most \fib{\textbf{1}}
    \end{avlpanel}
  \end{center}

  \note{The fix is simple to state. At every node, look at the height of the
  left side and the height of the right side. Do not let them differ by more
  than one. That is the whole idea.}
\end{frame}

% ------------------------------------------------------------------ DEFINITION
\begin{frame}[label=p1]{Definition: what an AVL tree is}
  \vspace{0.2em}
  \begin{columns}[T,onlytextwidth]
    \begin{column}{0.55\textwidth}
      \vspace{0.1em}
      {\avlfSec\color{avlmuted}Adelson-Velsky and Landis, 1962. The first
       self-balancing search tree.}\\[0.8em]

      \begin{avlpanel*}
        Balance factor of a node\\[6\pxl]
        $\mathrm{bf} = (\text{height of left}) - (\text{height of right})$
      \end{avlpanel*}\\[0.7em]

      {\strong{An AVL tree is:}}
      \begin{itemize}
        \setlength{\itemsep}{0.25em}
        \item a binary search tree, \emph{and}
        \item every node has \ok{$\mathrm{bf} = -1$, $0$, or $+1$}
      \end{itemize}
      \vspace{0.4em}
      {\avlfSec\color{avlmuted}If any node reaches $+2$ or $-2$, the tree is
       \viol{broken} and must be fixed.}
    \end{column}
    \begin{column}{0.42\textwidth}
      \begin{center}
      \begin{tikzpicture}[x=1cm,y=1cm,scale=0.92,every node/.style={transform shape}]
        \node[avlkeyok] (a) at (0,0) {};
        \node[avlannotok,anchor=west] at (0.40,0.16) {$+1$};
        \node[avlkey] (b) at (-1.15,-1.25) {};
        \node[avlannotok,anchor=east] at (-1.52,-1.06) {$0$};
        \node[avlkey] (c) at (1.15,-1.25) {};
        \node[avlannotok,anchor=west] at (1.52,-1.06) {$0$};
        \node[avlkey] (d) at (-1.9,-2.5) {};
        \node[avlkey] (e) at (-0.4,-2.5) {};
        \foreach \p/\q in {a/b,a/c,b/d,b/e} \draw[avledge] (\p) -- (\q);
        \node[avlannot] at (0,-3.30) {left height 2, right height 1};
        \node[color=avlbalance,font=\avlfSec] at (0,-3.85) {legal AVL tree};
      \end{tikzpicture}
      \end{center}
    \end{column}
  \end{columns}

  \note{Now the name and the rule. AVL stands for the two Soviet
  mathematicians who invented it in 1962. Balance factor is left height minus
  right height. An AVL tree is a BST where every single node has a balance
  factor of minus one, zero, or plus one.}
\end{frame}

% ------------------------------------------------------------------ how it solves it
\begin{frame}[label=p1]{How AVL fixes the bad case}
  \vspace{-0.2em}
  \begin{center}
  \begin{tikzpicture}[x=1cm,y=1cm,scale=0.86,every node/.style={transform shape}]
    \begin{scope}[xshift=-3.6cm,yshift=0.4cm]
      \foreach \k [count=\i from 0] in {1,2,3,4,5,6,7} {
        \pgfmathsetmacro{\yy}{-\i*0.55}
        \pgfmathsetmacro{\xx}{\i*0.30}
        \node[avlkeyalert,minimum size=36\pxl,font=\avlfAnnot] (c\k) at (\xx,\yy) {\k};
      }
      \foreach \a/\b in {c1/c2,c2/c3,c3/c4,c4/c5,c5/c6,c6/c7}
        \draw[draw=avlalert,line width=1.6\pxl] (\a) -- (\b);
      \node[avlannot,text=avlmuted] at (0.9,-4.2) {plain BST};
      \node[color=avlalert,font=\avlfLead\wbold] at (0.9,-4.85) {height 6};
    \end{scope}

    % neutral arrow: the transformation is not part of the Fibonacci thread
    \draw[avlarrow] (-0.9,-1.6) -- (0.5,-1.6);
    \node[avlannot,above=2pt] at (-0.2,-1.55) {same input};

    \begin{scope}[xshift=3.4cm,yshift=-0.6cm]
      \node[avlkeyok] (a4) at (0,0)        {4};
      \node[avlkeyok] (a2) at (-1.3,-1.15) {2};
      \node[avlkeyok] (a6) at (1.3,-1.15)  {6};
      \node[avlkeyok] (a1) at (-2.0,-2.3)  {1};
      \node[avlkeyok] (a3) at (-0.6,-2.3)  {3};
      \node[avlkeyok] (a5) at (0.6,-2.3)   {5};
      \node[avlkeyok] (a7) at (2.0,-2.3)   {7};
      \foreach \a/\b in {a4/a2,a4/a6,a2/a1,a2/a3,a6/a5,a6/a7}
        \draw[draw=avlbalance!70,line width=1.6\pxl] (\a) -- (\b);
      \node[avlannot,text=avlmuted] at (0,-3.2) {AVL tree};
      \node[color=avlbalance,font=\avlfLead\wbold] at (0,-3.85) {height 2};
    \end{scope}
  \end{tikzpicture}
  \end{center}

  \vspace{0.7em}
  \begin{center}
    {\avlfSec\color{avlmuted}Insert 1, 2, 3, 4, 5, 6, 7 into both.
     AVL rebalances as it goes. \ok{7 steps become 3.}}
  \end{center}
  \vspace{0.5em}

  \note{Same sorted input that destroyed the plain tree. The AVL tree fixes
  itself while you insert, and stays short. We will see exactly how it does
  that in Part 2.}
\end{frame}

% ------------------------------------------------------------------ height property
\begin{frame}[label=p1]{The height property}
  \vspace{0.1em}
  \begin{center}
    {\avlfLead\color{avlbody}A local rule at every node gives a
     \emph{global} guarantee:}\\[0.9em]
    \begin{avlpanel}[0.50\textwidth]
      {\avlfDivider\color{avlfib}$h < 1.44 \, \log_2 n$}
    \end{avlpanel}
  \end{center}

  % Three columns, a single hairline under the header, no rules between rows
  % and no fills: the payoff cell is carried by colour alone.
  \vspace{0.9em}
  \begin{center}
  \setlength{\tabcolsep}{20\pxl}
  \renewcommand{\arraystretch}{1.45}
  \setlength{\arrayrulewidth}{1\pxl}
  \arrayrulecolor{avlrule}
  {\avlfSec
  \begin{tabular}{lcc}
    \color{avlfaint}nodes & \color{avlfaint}perfect tree
      & \color{avlfaint}worst AVL tree\\\hline
    $1\,000\,000$ & 20 steps & \fib{\textbf{28 steps}}\\
  \end{tabular}}
  \end{center}

  \vspace{0.7em}
  \begin{center}
    {\avlfSec\color{avlmuted}Never a list. Always close to the best possible.}
  \end{center}

  \note{This is the promise. No matter what order the data arrives in, the
  height stays within 44 percent of a perfect tree. One million records: at
  most 28 comparisons. Next question --- why 1.44? Where does that number come
  from?}
\end{frame}

% ------------------------------------------------------------------ proof setup
%  The framing beat only. The recurrence itself is NOT stated here any more --
%  it is built from actual small trees on the next slide, so it arrives as
%  something the audience has already watched happen four times.
\begin{frame}[label=p1]{Proof: ask the question backwards}
  \vspace{0.9em}
  \begin{center}
    {\avlfLead\color{avlmuted}Instead of ``how tall can it get?'', ask:}\\[0.8em]
    % the pivot line: the turned question, with the turn itself in ink
    {\avlfStmt\color{avlfib}\wbold What is the
     {\color{avlink}\mdseries smallest} AVL tree of height $h$?}
  \end{center}

  \vspace{0.7em}
  \centerline{{\color{avlrule}\rule[2\pxl]{70\pxl}{1\pxl}}\hskip10\pxl
              {\color{avlfaint}\rule{4\pxl}{4\pxl}}\hskip10\pxl
              {\color{avlrule}\rule[2\pxl]{70\pxl}{1\pxl}}}

  \vspace{0.7em}
  \begin{center}
    {\avlfBody\color{avlbody}If even the emptiest legal tree of height $h$
     needs a lot of nodes,\\[0.3em]
     then a tree with few nodes \emph{cannot} be tall.}\\[0.9em]
    \begin{avlpanel}[0.58\textwidth]
      Write \fib{$N(h)$} for that smallest node count.
    \end{avlpanel}
  \end{center}

  \note{Turn the question around. Instead of bounding the tallest tree, build
  the emptiest legal tree of a given height and count it. If the emptiest tree
  of height h already needs a lot of nodes, then a tree with few nodes cannot
  possibly be that tall. Same fact, and this direction we can actually count.}
\end{frame}

% ------------------------------------------------------------------ BUILD-UP
%  The four smallest minimal trees, revealed one at a time. Each new tree is
%  visibly the two before it, joined under one root -- so the recurrence is
%  SEEN before it is written down.
%
%  Coordinate map (per panel, origin at the panel's root):
%    T_0  a(0,0)
%    T_1  a(0,0)  b(-0.45,-0.8)
%    T_2  a(0,0)  b(-0.62,-0.8) c(-1.05,-1.6)   d(0.62,-0.8)
%    T_3  a(0,0)  b(-0.85,-0.8) c(-1.35,-1.6) d(-1.75,-2.4) e(-0.42,-1.6)
%                 f(0.85,-0.8)  g(0.45,-1.6)
%  Panel origins: -5.0, -2.9, 0.0, 3.8. Captions on a common baseline.
%  NOTE: the style key must not be called "cap" -- that collides with TikZ's
%  own line-cap key and the frame will not compile.
% Title kept to a single line on purpose: a second line pushes the rule down
% and steals the height the ladder and its panel need.
\begin{frame}[label=p1]{Each tree is the two before it, under one root}
  \vspace{-0.7em}
  \begin{center}
  % Vertical axis only (y=0.86cm), and no `transform shape`: the ladder loses
  % height so the panel below gets air, while the horizontal pitch stays at
  % 1 cm -- the T_k / N(h) captions are set at full size and would collide with
  % each other if the columns closed up. Node discs shrink by their own key.
  \begin{tikzpicture}[x=1cm,y=0.86cm,
      mini/.style   = {circle, draw=avlfib, fill=avlfib!12!avlground,
                       line width=1.3\pxl, inner sep=0pt, minimum size=4.0mm},
      miniroot/.style={circle, draw=avlfib, fill=avlfib, inner sep=0pt,
                       line width=1.3\pxl, minimum size=4.0mm},
      medge/.style  = {avlminiedge},
      blk/.style    = {draw=avlfaint, dashed, line width=1\pxl,
                       rounded corners=2pt, inner sep=2.2pt},
      blab/.style   = {avlannot},
      treecap/.style= {font=\avlfSec, text=avlink, align=center,
                       minimum width=2cm, minimum height=8mm}]
    % fixed extent so nothing shifts between overlays
    % Extent trimmed to what the beats actually occupy (deepest caption sits at
    % -3.85), so the panel below gets its room without shrinking the figure.
    \path (-5.9,-3.95) rectangle (5.5,0.30);

    % ---------------- T_0 ----------------
    \begin{scope}[xshift=-5.0cm]
      \node[miniroot] (za) at (0,0) {};
      \node[treecap] at (0,-3.25) {$T_0$\\ \fib{$N(0)=1$}};
    \end{scope}

    % ---------------- T_1 ----------------
    \only<2->{\begin{scope}[xshift=-2.9cm]
      \node[miniroot] (oa) at (0,0) {};
      \node[mini] (ob) at (-0.45,-0.8) {};
      \draw[medge] (oa) -- (ob);
      \node[treecap] at (0,-3.25) {$T_1$\\ \fib{$N(1)=2$}};
    \end{scope}}

    % ---------------- T_2 = root + T_1 + T_0 ----------------
    \only<3->{\begin{scope}[xshift=0cm]
      \node[miniroot] (ta) at (0,0) {};
      \node[mini] (tb) at (-0.62,-0.8) {};
      \node[mini] (tc) at (-1.05,-1.6) {};
      \node[mini] (td) at ( 0.62,-0.8) {};
      \draw[medge] (ta) -- (tb);  \draw[medge] (tb) -- (tc);
      \draw[medge] (ta) -- (td);
      \node[blk, fit=(tb)(tc), label={[blab]below:$T_1$}] {};
      \node[blk, fit=(td),     label={[blab]below:$T_0$}] {};
      \node[treecap] at (0,-3.25) {$T_2$\\ \fib{$N(2)=1+2+1=4$}};
    \end{scope}}

    % ---------------- T_3 = root + T_2 + T_1 ----------------
    \only<4->{\begin{scope}[xshift=3.8cm]
      \node[miniroot] (ha) at (0,0) {};
      \node[mini] (hb) at (-0.85,-0.8) {};
      \node[mini] (hc) at (-1.35,-1.6) {};
      \node[mini] (hd) at (-1.75,-2.4) {};
      \node[mini] (he) at (-0.42,-1.6) {};
      \node[mini] (hf) at ( 0.85,-0.8) {};
      \node[mini] (hg) at ( 0.45,-1.6) {};
      \draw[medge] (ha) -- (hb);  \draw[medge] (hb) -- (hc);
      \draw[medge] (hc) -- (hd);  \draw[medge] (hb) -- (he);
      \draw[medge] (ha) -- (hf);  \draw[medge] (hf) -- (hg);
      \node[blk, fit=(hb)(hc)(hd)(he), label={[blab]below:$T_2$}] {};
      \node[blk, fit=(hf)(hg),         label={[blab]below:$T_1$}] {};
      \node[treecap] at (0,-3.45) {$T_3$\\ \fib{$N(3)=1+4+2=7$}};
    \end{scope}}
  \end{tikzpicture}
  \end{center}

  % The caption block reserves its own height, so no beat shifts the figure.
  \begin{center}
  \begin{minipage}[t][112\pxl][c]{0.86\textwidth}\centering
    \only<1>{{\avlfSec\color{avlmuted}The smallest legal tree of height 0 is
              one node.}}
    \only<2>{{\avlfSec\color{avlmuted}For height 1 you need a child --- and one
              child is enough, because a gap of 1 is legal.}}
    \only<3>{{\avlfSec\color{avlmuted}For height 2, one side must reach
              height 1. The other is as \emph{small} as the rule allows:
              height 0.}}
    \only<4>{{\avlfSec\color{avlmuted}And again: one side $T_2$, the other side
              $T_1$. \fib{The pattern does not change.}}}
    % the recurrence arrives in a panel -- it is the payoff, not a caption
    \only<5->{\begin{avlpanel}[0.74\textwidth]
                {\avlfTitle\color{avlfib}$N(h) = N(h-1) + N(h-2) + 1$}\\[8\pxl]
                {\avlfEyebrow\color{avlmuted}one root, plus the two smallest
                 trees that fit under it}
              \end{avlpanel}}
  \end{minipage}
  \end{center}

  \note{Click through them slowly. Height zero: one node. Height one: add a
  child. Height two: the left side has to reach height one, and the right side
  is as small as the rule permits, which is height zero. Point at the two
  dashed boxes. Height three: same again --- a T two and a T one. By now they
  should be saying the recurrence before you write it. Then click and write it.}
\end{frame}

% ------------------------------------------------------- GOLDEN RATIO (algebra)
\begin{frame}[label=p1]{Where the golden ratio comes from}
  % Beats 1-4 build the algebra top-down: a faint eyebrow announces each move,
  % the maths follows in serif, and the line that matters turns ember. Beat 5
  % clears the board for the payoff -- phi itself.
  \only<1-4>{
  \vspace{-0.8em}
  \begin{center}
    {\avlfSec\color{avlmuted}The $+1$ is in the way --- so absorb it.}\\[0.6em]
    {\avlfLead $N(h) = N(h-1) + N(h-2) + 1$}\\[0.6em]

    \onslide<2->{%
      \eyebrow{add 1 to both sides, and write $M(h) = N(h)+1$}\\[0.4em]
      {\avlfBody $\big(N(h)+1\big) = \big(N(h{-}1)+1\big) + \big(N(h{-}2)+1\big)$}}
    \\[0.7em]

    \onslide<3->{%
      {\avlfLead\color{avlfib}$M(h) = M(h-1) + M(h-2)$}\\[0.25em]
      {\avlfEyebrow\color{avlmuted}Fibonacci exactly, no $+1$ left}}
  \end{center}

  \onslide<4->{
  \vspace{0.1em}
  \begin{center}
    {\color{avlrule}\rule{0.72\textwidth}{1\pxl}}\\[0.3em]
    \eyebrow{look for solutions of the form $M(h) = x^h$}\\[0.25em]
    {\avlfBody $x^h = x^{h-1} + x^{h-2}
      \;\Longrightarrow\; x^2 = x + 1
      \;\Longrightarrow\; x = (1 \pm \sqrt5)\,/\,2$}
  \end{center}}}

  % Beat 5 -- the phi payoff. A split panel: the glyph, a rule, the value.
  \only<5>{
  \vspace{0.1em}
  \begin{center}
    \eyebrow{$x^2 = x+1$ has one root above 1, and it dominates}\\[1.0em]
    % Laid out in tikz so the glyph, the rule and the two value lines all
    % share one vertical centre -- baseline stacking put phi low in the box.
    % The glyph and its value, filling the panel: phi set large and bold in
    % ink, the number in ember so the two are told apart at a glance.
    \begin{avlpanel}[0.38\textwidth]
      \begin{tikzpicture}[baseline=(mid.base)]
        \coordinate (mid) at (0,-0.08);
        \node[anchor=east,inner sep=0pt] at (-0.30,0)
          {\fontsize{50}{50}\selectfont\color{avlink}$\boldsymbol{\varphi}$};
        \draw[avlrule,line width=1\pxl] (0,-0.46) -- (0,0.46);
        \node[anchor=west,inner sep=0pt] at (0.30,0)
          {\fontsize{26}{26}\selectfont\wbold\color{avlfib}1.618};
      \end{tikzpicture}
    \end{avlpanel}\\[0.9em]
    {\avlfTitle\color{avlink}$N(h) \sim \varphi^{\,h}$}\\[0.9em]
    \begin{minipage}{0.74\textwidth}\centering
      {\avlfBody\color{avlbody}The smallest legal tree multiplies by
       \fib{1.618 every level}. Nodes grow exponentially, so height can only
       grow like a logarithm --- \fib{this is the whole reason 1.44 appears in
       this talk.}}
    \end{minipage}
  \end{center}}

  \note{The recurrence is almost Fibonacci but for the plus one. Adding one to
  both sides makes it vanish --- M of h is exactly Fibonacci. Then the standard
  trick: guess a geometric solution, divide through by x to the h minus two,
  and you are left with x squared equals x plus one. That is the defining
  equation of the golden ratio. So the smallest AVL tree grows like phi to the
  h --- and that is the only reason 1.44 appears anywhere in this talk.}
\end{frame}

% --------------------------------------------------------- the exact identity
\begin{frame}[label=p1]{The exact count}
  % The identity is set bare, big and ember -- it needs no box; the panel at
  % the foot of the frame is where the consequence lands.
  \vspace{0.2em}
  \begin{center}
    {\avlfDivider\color{avlfib}$N(h) = F(h+3) - 1$}\\[0.7em]
    {\avlfSec\color{avlmuted}because $M(h) = N(h)+1$ is Fibonacci, starting at
     $M(0) = 2 = F(3)$}
  \end{center}

  % No vertical rules, no zebra: a hairline under the header, avlrule2 between
  % body rows, and the payoff column carries the ember tint.
  \vspace{0.7em}
  \onslide<2->{
  \begin{center}
  \setlength{\tabcolsep}{10\pxl}
  \renewcommand{\arraystretch}{1.35}
  \setlength{\arrayrulewidth}{1\pxl}
  {\avlfSec
  \begin{tabular}{r|cccccccc}
    \color{avlfaint}$h$ & 0 & 1 & 2 & 3 & 4 & 5 & 6 & \fib{\textbf{7}}\\
    \color{avlfaint}$N(h)$ & 1 & 2 & 4 & 7 & 12 & 20 & 33 & \fib{\textbf{54}}\\
    \color{avlfaint}$F(h{+}3)$ & 2 & 3 & 5 & 8 & 13 & 21 & 34 & \fib{\textbf{55}}\\
  \end{tabular}}
  \end{center}}

  \vspace{0.8em}
  \onslide<3->{
  \begin{center}
    \begin{avlpanel}[0.58\textwidth]
      \panellabel{taking logs of $N(h) \sim \varphi^h$}
      {\avlfTitle\color{avlfib}$h \approx \log_2 n \,/\, \log_2\varphi
        = 1.44 \log_2 n$}
    \end{avlpanel}
  \end{center}}

  \note{One line. The plus-one substitution we just made IS the minus one here.
  Read the last column out loud: at height seven the smallest AVL tree has 54
  nodes, and 55 is a Fibonacci number. Hold onto 54 --- Part 3 will build this
  exact tree and break it.}
\end{frame}

% ------------------------------------------------------------------ THE PLANT
\begin{frame}[label=p1]{Remember this number: 54}
  \vspace{-0.4em}
  \begin{center}
  % height-7 Fibonacci tree: leafSep 44 px, levSep 38 px (spec section 5)
  \begin{tikzpicture}[x=1cm,y=1cm,
      avldot/.style={circle,draw=avlfib,fill=avlfib,inner sep=0pt,
                     minimum size=5.8\pxl},
      avledge/.style={draw=avlfib!75,line width=1.3\pxl}]
    \setavlscale{0.48}{0.37}
    \resetavl
    \drawfib{7}{0}{0}
  \end{tikzpicture}
  \end{center}

  \vspace{0.1em}
  \begin{center}
    {\avlfLead\color{avlink}This is the \fib{worst tree AVL allows} at
     height 7.}\\[0.4em]
    {\avlfSec\color{avlmuted}It has the fewest nodes possible:
     $N(7) = F(10) - 1 = 55 - 1 =$ \fib{\textbf{54}}}\\[0.7em]
    % the Part 3 plant, in Part 3's hue -- the only plum on a Part 1 slide
    \begin{avlpanelh}[0.98\textwidth]{avlpthree}
      {\color{avlpthree}In Part 3 we will delete one leaf from this exact tree
       --- and watch it come apart.}
    \end{avlpanelh}
  \end{center}

  \note{This is the emptiest legal AVL tree of height seven. Fibonacci predicts
  it has 54 nodes. Remember that number. In Part 3 we will take one leaf out of
  it and count the damage. Handoff: we know the tree stays short. Now, how does
  it fix itself?}
\end{frame}
